\setcounter{secnumdepth}{3}
\section{Вступ}

\subsection{Актуальність}

Автоматичний переклад мовлення є однією з найскладніших задач на перетині
обробки природної мови, акустичного моделювання та машинного перекладу.
Розвиток нейромережевих технологій, зокрема архітектури
Transformer~\cite{vaswani}, призвів до появи каскадних та end-to-end
систем~\cite{cascade_st, st_survey}, здатних здійснювати переклад
мовлення. Проте питання якості такого перекладу залишається відкритим,
особливо для мовних пар, де одна з мов є морфологічно багатою та менш
представленою у навчальних даних~\cite{low_resource_mt}.

Пара українська--англійська становить інтерес з кількох причин. Українська мова
має складну морфологічну систему та відносно вільний порядок слів, що створює
додаткові виклики для автоматичного розпізнавання та перекладу. Обсяг доступних
паралельних корпусів для цієї пари значно менший, ніж для поширених мовних
пар на кшталт англійська--французька чи англійська--німецька. В умовах активної
міжнародної комунікації України потреба в якісному автоматичному перекладі
з української на англійську зростає, зокрема у сфері політичного дискурсу.

Більшість досліджень якості перекладу мовлення зосереджені на оцінюванні
окремих компонентів системи --- точності розпізнавання мовлення або якості
машинного перекладу --- без аналізу того, як помилки накопичуються вздовж
усього конвеєра~\cite{error_propagation, bentivogli_cascade} та як вони
впливають на зв'язність перекладу для кінцевого слухача. Комплексне
дослідження семантичної точності, каскадного ефекту та міжфразової
когерентності перекладу саме для мовної пари українська--англійська
визначає актуальність даної роботи.

\subsection{Об'єкт, предмет, мета і завдання дослідження}

\textit{Об'єктом} дослідження є якість україно-англійського автоматичного
перекладу усного мовлення.

\textit{Предметом} дослідження є семантична точність та міжфразова когерентність
перекладу, здійснюваного каскадними системами автоматичного перекладу мовлення
(ASR + MT).

\textit{Метою} кваліфікаційної роботи є кількісне оцінювання якості
україно-англійського автоматичного перекладу усного мовлення з акцентом на
семантичну точність, каскадний ефект поширення помилок та міжфразову
когерентність.

Досягнення поставленої мети вимагає розв'язання таких задач:
\begin{enumerate}[nosep]
    \item Розглянути сучасні підходи до автоматичного перекладу мовлення
    та методи оцінювання його якості.
    \item Визначити особливості мовної пари українська--англійська, що впливають
    на якість автоматичного перекладу.
    \item Сформувати корпус дослідження на основі записів політичних звернень
    українською мовою та відповідних референсних перекладів.
    \item Обрати та обґрунтувати систему метрик для оцінювання семантичної
    точності, когерентності та інших характеристик перекладу.
    \item Провести експериментальне оцінювання якості перекладу обраних систем
    (спеціалізованої NMT та LLM-моделей) та виконати статистичний аналіз
    результатів.
    \item Дослідити каскадний ефект --- вплив помилок розпізнавання мовлення
    на якість машинного перекладу.
    \item Порівняти якість перекладу спеціалізованої NMT-моделі та великих
    мовних моделей (LLM) для цільової мовної пари та домену.
\end{enumerate}

\subsection{Методи дослідження}

У дослідженні застосовано кількісні та якісні методи оцінювання перекладу.
Для автоматичного оцінювання семантичної точності використано метрики
BLEU~\cite{papineni_bleu}, BERTScore~\cite{bert_score}
та METEOR~\cite{banerjee_meteor}. Для оцінювання міжфразової когерентності розроблено
методику аналізу контекстуальної зв'язності між послідовними сегментами
перекладу. Додатково застосовано класифікацію помилок перекладу за типами:
пропуски, спотворення та додавання.

Для статистичного аналізу результатів використано кореляційний аналіз
(коефіцієнти кореляції Пірсона та Спірмена), критерій Манна-Уітні для
порівняння MT-систем та критерій Краскела-Уолліса для множинного
порівняння LLM-моделей.

\subsection{Наукова новизна}

Наукова новизна роботи полягає у:
\begin{itemize}[nosep]
    \item комплексному кількісному оцінюванні якості україно-англійського
    автоматичного перекладу усного мовлення з використанням системи
    автоматичних метрик (BLEU, BERTScore, METEOR) та ручної анотації помилок;
    \item кількісному дослідженні каскадного ефекту поширення помилок ASR
    на якість MT з виявленням статистично значущих кореляцій;
    \item порівняльному аналізі спеціалізованої NMT-моделі та LLM-систем
    для перекладу політичного дискурсу з української на англійську;
    \item аналізі міжфразової когерентності як окремого критерію якості
    перекладу для мовної пари українська--англійська.
\end{itemize}

\subsection{Практичне значення}

Результати роботи можуть бути використані для вдосконалення систем
автоматичного перекладу мовлення, зокрема для ідентифікації критичних
категорій помилок (власні назви, географічні назви, термінологія).
Запропонована методологія оцінювання може застосовуватися для порівняльного
аналізу різних систем перекладу, а також для оцінювання якості перекладу
у сферах дипломатичної та політичної комунікації. Кількісне порівняння
чотирьох MT-систем надає практичні рекомендації щодо вибору компонентів
каскадного pipeline залежно від вимог до латентності та якості.

\subsection{Апробація}

Окремі результати дослідження були представлені на студентських семінарах та
під час внутрішніх обговорень на кафедрі.

\clearpage
